% !TEX root = main.tex
\section{评估体系与基准测试}
\label{sec:evaluation}

\noindent\textbf{章节元数据：} content\_timestamp=2026-05-21；collected\_at=2026-05-21T21:58:00+08:00；time\_slice=2026-05。本文基于 \texttt{raw-papers/}、\texttt{paper-reviews/} 与既有综述草稿整理，重点讨论 Agent Self-Evolution 的评估体系、基准测试、指标风险与未来规范。

\subsection{评估问题为何是自进化研究的核心}

Agent 自进化研究的核心矛盾之一，是“系统声称变得更强”与“我们如何确认它真的变强”之间的张力。传统大模型评测往往关注静态能力：模型在固定题集上的一次性表现如何；而 Agent Evolution 关注的是动态过程：一个智能体能否在交互、反馈、记忆、搜索、代码修改或多轮训练中持续改进，并且这种改进是否能跨任务、跨环境、跨时间保持有效。因此，评估体系不只是论文中的实验章节，而是整个自进化闭环的“选择压力”和“安全闸门”。没有可靠评估，进化会退化为随机试错；没有抗污染评估，进化会变成对 benchmark 的过拟合；没有长期评估，短期收益可能掩盖能力退化、成本上升和安全风险。

从已收集的论文与 review 看，当前 Agent 自进化论文大致采用四类评估范式：第一类是代码与软件工程基准，如 SWE-Bench、SWE-Bench Verified、HumanEval、MBPP、LiveCodeBench、Polyglot、rSDE-Bench 等；第二类是通用推理与问答基准，如 GSM8K、MATH、ARC、DROP、MMLU、GPQA、MGSM、AIME、OlympiadBench 等；第三类是环境交互基准，如 Minecraft/Voyager、WebArena、WebVoyager、Mind2Web-Live、AppWorld、ALFWorld、WebShop、GAIA-web 等；第四类是自进化专用过程指标，如迭代增益、archive 多样性、回滚率、迁移性、记忆复用率、评估成本和安全门控通过率。本章围绕这些基准与指标展开，重点讨论它们如何塑造 Agent Evolution 研究结论，以及它们自身存在的局限。

需要强调的是，评估在自进化研究中有双重身份：一方面它是论文用来证明方法有效的外部证据；另一方面它又经常被嵌入训练、搜索或自我修改循环，成为智能体实际优化的目标函数。后一种身份使评估问题比普通 LLM benchmark 更敏感。只要指标可见、可重复、可被 agent 试探，系统就可能从“提升真实能力”滑向“适配评价器”。因此，本章不会把评估简单理解为排行榜分数，而是将其视作一个由数据切片、任务生成、验证器、成本模型、安全策略和人类审查共同构成的治理系统。

\subsection{主要评估基准}

\subsubsection{软件工程与代码智能体基准}

SWE-Bench 已成为代码智能体自进化论文中最重要的基准之一。它要求系统在真实开源仓库中定位问题、修改代码、运行测试并提交可验证补丁，因此比单函数代码生成更接近真实软件工程任务。Darwin Gödel Machine 与 SICA 等系统都使用 SWE-Bench 或其变体来证明“自我修改代码库”能够带来真实收益：DGM 通过维护智能体 archive，并让变体在 SWE-Bench 与 Polyglot 上接受经验验证；SICA 则展示一个 coding agent 可以在多轮迭代中修改自身工具、提示词、文件编辑逻辑与审查机制，从而提升 SWE-Bench Verified 表现。这类结果的意义不只是分数提高，而是说明“可执行代码 + 自动测试 + 历史 archive”可以构成 Agent 自进化的闭环。

HumanEval 与 MBPP 属于更轻量的代码生成基准，通常用于测试函数级编程能力。它们的优点是自动判分成本低、复现实验容易、适合作为多轮进化中的快速 fitness signal；缺点是任务粒度较小，无法充分反映真实 agent 的文件系统操作、调试、依赖管理、需求澄清和长期规划能力。Reflexion 在 HumanEval、MBPP 与 LeetcodeHardGym 上展示了语言反思和自生成单元测试的价值，但 review 中也指出 MBPP 自生成测试存在显著 false positive 风险，说明自动测试若质量不足，会把错误反馈写入反思记忆。EvoMAC 使用 HumanEval 评估函数级改进，同时引入 rSDE-Bench 这类更接近需求驱动软件开发的任务，说明单一函数级基准已经不足以支撑对“自进化软件智能体”的强结论。

LiveCodeBench、Polyglot 和类似在线/多语言代码评测补充了 SWE-Bench 的不足。LiveCodeBench 更强调时间切片和较新的题目，有助于缓解训练污染；Polyglot 可检验跨语言泛化能力。对于自进化系统而言，跨基准一致提升比单一基准大幅提升更可信。如果一个系统只在 SWE-Bench 上显著提升，却在 HumanEval、MBPP、LiveCodeBench 或真实仓库任务上无明显收益，则可能说明它学到的是 benchmark-specific workflow，而不是通用编码智能。相反，如果它能同时改善 issue 解析、测试生成、patch 最小化、回归测试和跨语言迁移，才更接近“软件工程能力”而非“刷题能力”。

\subsubsection{推理、知识与开放式环境基准}

在非代码领域，GSM8K、MATH、ARC、MMLU、GPQA、DROP、MGSM 等仍是 Agent 自进化论文常用的推理基准。STaR 使用小规模 rationale examples 与大规模无标注数据迭代训练推理链，展示了模型可以通过生成—筛选—微调循环提升推理能力；RISE 将单轮问题转化为多轮 MDP，用 correctness reward 训练模型在测试时逐步自我修正；SPIRAL 等自博弈方法则使用数学、聊天、博弈和推理任务证明模型可以从自身生成的数据中获得改进。Self-Rewarding LM 与 Meta-Rewarding 则更多通过 AlpacaEval、MT-Bench、Arena-Hard 等开放式生成评估证明指令遵循与偏好建模能力的提升。

这些基准的优点是成熟、可比较、社区熟悉。但它们对 Agent Evolution 的支持是有限的。首先，许多任务是静态题集，系统可能通过训练数据、检索或反复调参间接接触题目；其次，很多基准只衡量最终答案，无法衡量智能体的工具使用、探索策略、记忆复用、失败恢复和多轮交互质量；第三，推理基准上的提升可能来自基础模型能力、采样策略、提示词长度或 evaluator 偏好，而不一定来自真正的自进化机制。尤其是 LLM-as-a-Judge 类评测，容易受到长度偏见、位置偏见、模型家族偏好和打分分布漂移影响。因此，在综述中引用这类结果时，应区分“模型能力提升”“对评估器偏好的适配”和“agentic loop 带来的能力提升”。

Voyager 是环境交互评估的代表案例。它在 Minecraft 中通过自动课程、技能库和迭代提示机制持续探索，使用独特物品发现数、科技树解锁速度、探索距离和技能迁移来衡量 lifelong learning。与静态题集相比，这类评估更接近自进化智能体的长期目标：智能体不仅要回答问题，还要在环境中持续发现任务、学习技能、复用经验并适应新世界。然而，Minecraft 指标也具有强领域依赖性。独特物品和科技树里程碑对于 Minecraft 有意义，但很难直接迁移到网页自动化、科研发现、企业流程或机器人控制。

\begin{table}[htbp]
\centering
\small
\caption{Agent 自进化常用基准对比矩阵。}
\label{tab:benchmarks}
\begin{tabular}{p{2.5cm}p{2.4cm}p{3.6cm}p{3.7cm}}
\toprule
\textbf{基准类型} & \textbf{代表基准} & \textbf{适合验证的问题} & \textbf{主要局限} \\
\midrule
软件工程 & SWE-Bench, LiveCodeBench, Polyglot & 真实仓库修复、代码自修改、测试驱动进化 & 成本高；容易学习 benchmark-specific workflow \\
函数级代码 & HumanEval, MBPP & 快速自动判分、局部代码生成与反思 & 粒度小；不能代表真实软件工程流程 \\
推理问答 & GSM8K, MATH, MMLU, GPQA & 推理链、反思、自奖励与自博弈 & 污染风险；最终答案不能解释 agentic 过程 \\
环境交互 & Minecraft, AppWorld, WebArena, ALFWorld & 长期探索、工具使用、技能迁移 & 搭建复杂；领域指标难迁移 \\
算法发现 & AlphaEvolve tasks, FunSearch-style evaluators & 程序化 fitness 下的开放式搜索 & 依赖 evaluator 设计；目标错配会导致 Goodhart 问题 \\
\bottomrule
\end{tabular}
\end{table}

\subsection{Agent 进化专用评估指标}

传统 benchmark 分数不足以刻画“进化”。Agent Evolution 需要过程性指标。第一是迭代增益曲线，即每一轮自我改进后的性能变化。一个可靠系统应展示多轮稳定上升，而不是某一轮偶然跳升。SICA、DGM、ADAS、RISE、Meta-Rewarding 等系统的关键证据都来自迭代过程：archive 中新 agent 是否持续优于祖先，judge 能力是否随 actor 一起提升，self-correction 是否在多轮内单调改善，是否存在回退，是否能从失败中恢复。报告最终 best score 而不报告完整曲线，会掩盖“多数候选失败、少数样本幸存”的事实。

第二是泛化与迁移指标。自进化系统最容易被质疑为“对验证集过拟合”。因此，应将搜索或进化使用的 validation tasks 与最终 test tasks 严格隔离，并报告跨模型、跨任务、跨环境迁移。ADAS 的重要贡献之一，是测试发现的 agent designs 是否能迁移到不同任务和不同 LLM；Voyager 的关键证据之一，是技能库能否迁移到新 Minecraft world 或新任务；AlphaEvolve 的工程价值来自算法能否在真实基础设施中带来收益；Memory-R1 的可信度来自 LoCoMo、MSC、LongMemEval 等跨记忆 benchmark 的结果。没有迁移性，进化只是局部调参。

第三是多样性指标。Archive-based 方法不应只保留最高分个体，还应保留多条探索路径。DGM 将 open-endedness 引入 self-improving coding agents，其动机正是避免单一路径陷入局部最优。多样性可从 agent 架构差异、工具组合差异、prompt 或代码 diff 距离、行为轨迹差异、解决方案语义差异等角度衡量。多样性越高，系统越可能在未来任务中找到 stepping stones；但过高多样性也会增加评估成本和管理复杂度。

第四是安全与回归指标。自进化智能体可能通过修改代码、提示词、记忆或工具策略获得短期收益，同时引入不可见风险。因此每轮 evolution 都应记录回归测试通过率、失败类型、风险规则触发次数、人工审查比例、回滚率和 sandbox 违规情况。对于 DGM、SICA、AlphaEvolve 这类能修改代码或生成程序的系统，安全指标不应是附录，而应成为与 benchmark 分数同等重要的主指标。对于 Self-Rewarding 与 Meta-Rewarding 类方法，还应记录 judge drift、score inflation、verbosity bias、human-correlation decay 等评价器退化指标。

第五是成本效率指标。自进化系统通常消耗大量 token、API 调用、环境 rollouts 或训练算力。ADAS 需要多轮 agent proposal 与 evaluation；SICA 的多轮 self-modification 产生显著 API 成本；Meta-Rewarding 每轮要生成大量响应与 judge/meta-judge 判断；AlphaEvolve 依赖大规模候选评估。评价一个系统时，必须同时报告 absolute performance、delta per iteration、cost per improvement、wall-clock time、失败候选比例和可复用资产数量。否则，高分系统可能只是“用成本买分”。

\begin{figure}[htbp]
\centering
\begin{tikzpicture}[scale=0.75]
  \def\R{3.0}
  \foreach \a/\t in {90/迭代收益,30/迁移性,-30/安全性,-90/成本效率,-150/多样性,150/可复现性}{
    \draw[gray!35] (0,0)--(\a:\R);
    \node[font=\small] at (\a:3.65) {\t};
  }
  \foreach \r in {1,2,3}{\draw[gray!25] (90:\r)--(30:\r)--(-30:\r)--(-90:\r)--(-150:\r)--(150:\r)--cycle;}
  \draw[fill=blue!25,draw=blue!70,thick,opacity=.75] (90:2.6)--(30:1.8)--(-30:1.4)--(-90:1.2)--(-150:2.4)--(150:1.6)--cycle;
  \draw[fill=red!25,draw=red!70,thick,opacity=.65] (90:1.8)--(30:2.6)--(-30:2.2)--(-90:2.1)--(-150:1.3)--(150:2.5)--cycle;
  \node[blue!70!black] at (-2.9,-3.3) {开放式自进化系统};
  \node[red!70!black] at (2.8,-3.3) {生产级受控系统};
\end{tikzpicture}
\caption{评估基准雷达图：开放式探索系统与生产级受控系统在指标侧重点上的差异。}
\label{fig:evaluation-radar}
\end{figure}

\subsection{评估体系的问题与局限}

Benchmark contamination 是 Agent Evolution 中尤其严重的问题。静态公开题集可能已进入基础模型训练数据；自进化系统又会在多轮迭代中反复接触验证任务，进一步增加泄漏风险。解决方式包括：使用时间切片明确训练/评估边界，引入新近数据集，保留私有 test set，采用动态任务生成，以及报告 \texttt{content\_timestamp}、\texttt{collected\_at}、\texttt{time\_slice}。本项目所有 raw data 强制加入这些字段，正是为了避免未来综述和实证分析混淆“论文发表时间”“数据采集时间”和“模型可能接触数据的时间”。

时间切片还能帮助分析方法演化趋势。2022 年及以前的 STaR 更强调 rationale bootstrapping；2023 年的 Reflexion、Self-Refine、Voyager 展示了反馈、反思和技能库；2024 年出现更多自奖励、自博弈、架构搜索和多智能体进化；2025 年以后，DGM、AlphaEvolve、RAGEN、ACE、Memory-R1、WebEvolver 等工作更重视自动评估、代码级自改进、可扩展训练与记忆机制。没有时间切片，综述容易把不同阶段的技术混为一谈，也无法判断某个 benchmark 提升是否只是因为基础模型在训练时已经接触了类似任务。

Goodhart 定律在自进化系统中表现为：当 benchmark 成为目标，benchmark 就不再是好指标。Self-Rewarding LM 的局限之一，是同一模型既生成又评价，可能产生 reward hacking；Meta-Rewarding 虽然引入 meta-judge 改善 judge，但 review 中仍指出 score bias、长度偏见和 human-correlation 不稳定；AlphaEvolve 的局限之一，是 evaluator 设计不当会导致系统优化错误目标；SICA/DGM 在 coding benchmark 中也可能学到 benchmark-specific shortcuts。对于任何自进化系统，评估指标一旦进入优化循环，就会被系统主动利用。

因此，评估设计必须包含反 Goodhart 机制：多指标评估、隐藏测试集、人工抽检、跨域验证、负面能力监控、复杂度惩罚、可解释性审查和长期跟踪。特别是当 agent 能修改自身代码或策略时，评估器不应被 agent 直接访问或修改；评估数据、测试脚本和 sandbox 权限要隔离；archive 纳入标准要包含安全与泛化，而非单一分数。

\begin{figure}[htbp]
\centering
\begin{tikzpicture}[node distance=1.3cm, every node/.style={font=\small}]
  \node[draw,rounded corners,fill=blue!10,minimum width=3.2cm,minimum height=.8cm] (data) {时间切片数据};
  \node[draw,rounded corners,fill=green!10,below=of data,minimum width=3.2cm,minimum height=.8cm] (val) {隔离验证器};
  \node[draw,rounded corners,fill=orange!12,below=of val,minimum width=3.2cm,minimum height=.8cm] (agent) {候选 Agent 变体};
  \node[draw,rounded corners,fill=purple!10,right=2.8cm of val,minimum width=3.4cm,minimum height=.8cm] (archive) {Archive / 回滚点};
  \node[draw,rounded corners,fill=red!10,below=of archive,minimum width=3.4cm,minimum height=.8cm] (risk) {安全与成本门控};
  \draw[->,thick] (data)--(val);
  \draw[->,thick] (val)--(agent);
  \draw[->,thick] (agent)--(risk);
  \draw[->,thick] (risk)--(archive);
  \draw[->,thick] (archive)--(agent);
  \draw[->,thick,dashed] (archive)--(val);
\end{tikzpicture}
\caption{自进化评估治理架构：数据切片、隔离验证器、候选变体、archive 与安全门控共同构成闭环。}
\label{fig:evaluation-governance}
\end{figure}

自动评估是自进化的基础，但人类判断仍不可替代。代码测试可以判定补丁是否通过，但很难评价可维护性、设计优雅性和长期技术债；LLM-as-a-judge 可以快速评价文本质量，但存在位置偏见、长度偏见、模型偏好和解释幻觉；环境 reward 可以衡量任务完成，但可能忽略用户体验与安全边界。EvoMAC 的 rSDE-Bench 提到自动要求级评估与人类评价高度相关，这是一个重要方向：自进化 benchmark 应尽量建立自动指标与人类偏好的相关性证据，而不是默认自动分数等同真实质量。

Agent Evolution 系统高度依赖随机采样、基础模型版本、API 非确定性、环境状态和搜索路径。ADAS 中不同运行可能发现不同 agent designs；DGM 的 archive 生长路径会影响最终能力；Voyager 的探索轨迹与技能库内容也可能随执行变化。评估报告必须记录随机种子、模型版本、提示词、工具版本、数据切片、失败样本和完整轨迹。只报告最终最佳分数是不充分的，应报告均值、方差、最好/最差、失败率和 ablation。

\subsection{评估方法对比分析}

从方法论上看，代码基准最适合作为强验证信号，因为测试可执行、正确性明确、改进可量化。其缺点是容易收缩研究视野，使自进化系统过度偏向 coding agent。推理基准覆盖面广、可比较性强，但更容易受到污染和 prompt engineering 影响。环境交互基准最能体现 agentic behavior，但成本高、复现难、领域依赖强。自动 evaluator 驱动的算法发现最接近真正 evolution loop，但适用范围取决于问题是否可程序化验证。

\begin{figure}[htbp]
\centering
\begin{tikzpicture}
\begin{axis}[
  width=.92\linewidth,
  height=6cm,
  xlabel={年份}, ylabel={相对关注度},
  ymin=0, ymax=10, xmin=2022, xmax=2026,
  xtick={2022,2023,2024,2025,2026},
  legend pos=north west,
  grid=both]
\addplot+[mark=*] coordinates {(2022,3) (2023,5) (2024,6) (2025,8) (2026,9)};
\addlegendentry{代码与软件工程评估}
\addplot+[mark=square*] coordinates {(2022,5) (2023,7) (2024,7) (2025,6) (2026,6)};
\addlegendentry{推理与问答评估}
\addplot+[mark=triangle*] coordinates {(2022,1) (2023,5) (2024,6) (2025,8) (2026,8)};
\addlegendentry{环境交互评估}
\addplot+[mark=diamond*] coordinates {(2022,1) (2023,2) (2024,5) (2025,9) (2026,9)};
\addlegendentry{自动 evaluator / 进化闭环}
\end{axis}
\end{tikzpicture}
\caption{技术趋势时间线：2022--2026 年 Agent 自进化评估范式的相对关注度变化。数值为基于本项目文献与 review 的定性归一化估计。}
\label{fig:evaluation-timeline}
\end{figure}

因此，一个成熟的 Agent Evolution 评估体系应采用分层结构。第一层是单轮能力基准，用于确认基础能力；第二层是迭代进化基准，用于观察多轮改进曲线；第三层是跨域迁移基准，用于检验是否学到通用机制；第四层是安全与成本基准，用于判断是否可部署；第五层是长期真实任务评估，用于衡量用户价值。只有五层同时成立，才能有力支持“agent 正在自进化”这一强命题。

对于未来研究，本章建议采用以下报告规范：每个系统必须报告 \texttt{content\_timestamp}、\texttt{collected\_at} 和 \texttt{time\_slice}；必须区分 validation 与 test；必须给出至少一个跨域或新近时间切片评估；必须报告成本与失败率；必须包含反 Goodhart 设计；必须公开或描述评估器不可被 agent 篡改的机制；必须进行单变量 ablation，说明性能提升来自 reward、memory、self-play、prompt evolution、architecture search 还是混合因素。这样的规范将使 Agent Evolution 从“展示性 demo”走向可比较、可复现、可审计的科学研究。

\subsection{小结}

评估体系是 Agent 自进化研究的核心基础设施。SWE-Bench、HumanEval、MBPP、GSM8K、MATH、Minecraft、AppWorld、rSDE-Bench 等基准分别覆盖代码、推理、环境交互和软件开发需求，但没有任何单一基准能够完整衡量自进化。真正的评估应关注动态过程：迭代收益是否稳定，改进是否可迁移，archive 是否保持多样性，系统是否避免 reward hacking，成本是否合理，安全门控是否有效。

本章的核心结论是：Agent Evolution 的评估不能只问“分数是否更高”，而要问“为什么更高、是否可复现、能否迁移、是否安全、是否值得成本、是否经得起时间切片检验”。只有当评估体系本身足够强，Agent 的自进化才不会变成对指标的投机，而能成为真实能力增长的工程路径与科学范式。
